% Abstract draft v2 (2026-06-11). Reflects the current energy-aware framing.
% Headline numbers are deliberately soft: the method is being reworked and OWNER
% is still running. Rewrite the last sentences once the final runs are complete.

Open-world named entity recognition (NER) must assign entity types that were not fixed at training time.
Recent systems either train a dedicated encoder or prompt billion-parameter language models, and they are compared almost exclusively on accuracy, ignoring their compute and energy cost.
We present \textbf{OPENER}, a lightweight open-world NER pipeline assembled from off-the-shelf components: a frozen mention detector, a Matryoshka-truncatable sentence embedder sharpened by a light contrastive fine-tuning, and a balanced linear classifier for entity typing, so that no encoder is trained from scratch.
We benchmark OPENER against GLiNER, GNER and OWNER on 13 datasets along three axes at once: quality, measured by Adjusted Mutual Information (AMI), inference latency, and energy.
Our analysis separates entity typing on gold mentions, where the contrastive embedding space is highly discriminative, from the end-to-end setting, where the dominant open-world bottleneck is shown to be mention detection rather than typing.
The energy-aware comparison further exposes order-of-magnitude efficiency gaps between approaches that a quality-only evaluation would hide.
\emph{[Headline results to be finalised once all runs, including OWNER, are complete.]}
We release the code, the per-experiment configurations and the full benchmark.
